% ============================================================================
% Referencias (autor-año). Las fuentes de la propuesta original conservan su enlace.
% ============================================================================
\renewcommand{\refname}{Referencias}
\begin{thebibliography}{99}
\small
\setlength{\itemsep}{2pt}

\bibitem[Adams y MacKay(2007)]{adams2007}
Adams, R.~P. y MacKay, D.~J.~C. (2007). \emph{Bayesian online changepoint detection}. arXiv:0710.3742.
\url{https://arxiv.org/abs/0710.3742}

\bibitem[A{\"i}t-Sahalia y Lo(1998)]{aitsahalia1998}
A{\"i}t-Sahalia, Y. y Lo, A.~W. (1998). Nonparametric estimation of state-price densities implicit in
financial asset prices. \emph{Journal of Finance}, 53(2), 499--547.

\bibitem[Alpaca(s.f.)]{alpaca}
Alpaca. \emph{Historical option data} (documentación). \url{https://docs.alpaca.markets/us/docs/historical-option-data}

\bibitem[Bailey y L{\'o}pez~de~Prado(2014)]{bailey2014}
Bailey, D.~H. y López de Prado, M. (2014). The deflated Sharpe ratio: correcting for selection bias,
backtest overfitting and non-normality. \emph{Journal of Portfolio Management}, 40(5), 94--107.
\url{https://www.davidhbailey.com/dhbpapers/deflated-sharpe.pdf}

\bibitem[Borovi{\v{c}}ka et~al.(2016)]{borovicka2016}
Borovička, J., Hansen, L.~P. y Scheinkman, J.~A. (2016). Misspecified recovery. \emph{Journal of
Finance}, 71(6), 2493--2544.

\bibitem[Boyd et~al.(2017)]{boyd2017}
Boyd, S., Busseti, E., Diamond, S., Kahn, R.~N., Koh, K., Nystrup, P. y Speth, J. (2017). Multi-period
trading via convex optimization. \emph{Foundations and Trends in Optimization}, 3(1), 1--76.
\url{https://stanford.edu/~boyd/papers/cvx_portfolio.html}

\bibitem[Breeden y Litzenberger(1978)]{breeden1978}
Breeden, D.~T. y Litzenberger, R.~H. (1978). Prices of state-contingent claims implicit in option prices.
\emph{Journal of Business}, 51(4), 621--651.

\bibitem[Burkovska et~al.(2018)]{burkovska2018}
Burkovska, O. et~al. (2018). Calibration to American options: numerical investigation of the
de-Americanization method. \emph{Quantitative Finance}. \url{https://pmc.ncbi.nlm.nih.gov/articles/PMC6034575/}

\bibitem[Carr y Wu(2009)]{carrwu2009}
Carr, P. y Wu, L. (2009). Variance risk premiums. \emph{Review of Financial Studies}, 22(3), 1311--1341.

\bibitem[Cboe DataShop(s.f.)]{cboe}
Cboe DataShop. \emph{Option quote intervals} (ficha de producto). \url{https://datashop.cboe.com/option-quote-intervals}

\bibitem[Chernozhukov et~al.(2010)]{chernozhukov2010}
Chernozhukov, V., Fernández-Val, I. y Galichon, A. (2010). Quantile and probability curves without
crossing. \emph{Econometrica}, 78(3), 1093--1125.

\bibitem[Cohen et~al.(2021)]{cohen2021}
Cohen, S.~N., Reisinger, C. y Wang, S. (2021). \emph{Arbitrage-free neural-SDE market models}.
arXiv:2105.11053. \url{https://arxiv.org/abs/2105.11053}

\bibitem[Cox et~al.(1979)]{cox1979}
Cox, J.~C., Ross, S.~A. y Rubinstein, M. (1979). Option pricing: a simplified approach. \emph{Journal of
Financial Economics}, 7(3), 229--263.

\bibitem[Demarta y McNeil(2005)]{demarta2005}
Demarta, S. y McNeil, A.~J. (2005). The t copula and related copulas. \emph{International Statistical
Review}, 73(1), 111--129.

\bibitem[Duan et~al.(2020)]{duan2020}
Duan, T., Avati, A., Ding, D.~Y., Thai, K.~K., Basu, S., Ng, A. y Schuler, A. (2020). NGBoost: natural
gradient boosting for probabilistic prediction. \emph{Proceedings of the 37th International Conference on
Machine Learning}. arXiv:1910.03225. \url{https://arxiv.org/abs/1910.03225}

\bibitem[Gatheral y Jacquier(2014)]{gatheral2014}
Gatheral, J. y Jacquier, A. (2014). Arbitrage-free SVI volatility surfaces. \emph{Quantitative Finance},
14(1), 59--71. arXiv:1204.0646. \url{https://arxiv.org/abs/1204.0646}

\bibitem[Gibbs y Cand{\`e}s(2021)]{gibbs2021}
Gibbs, I. y Candès, E. (2021). Adaptive conformal inference under distribution shift. \emph{Advances in
Neural Information Processing Systems}, 34. arXiv:2106.00170. \url{https://arxiv.org/abs/2106.00170}

\bibitem[GitHub(2026)]{github2026}
GitHub (2026). \emph{Ubuntu 26 generally available and latest migration}. GitHub Changelog, 17 de
septiembre de 2026.
\url{https://github.blog/changelog/2026-09-17-ubuntu-26-generally-available-and-latest-migration/}

\bibitem[Gneiting y Raftery(2007)]{gneiting2007}
Gneiting, T. y Raftery, A.~E. (2007). Strictly proper scoring rules, prediction, and estimation.
\emph{Journal of the American Statistical Association}, 102(477), 359--378.
\url{https://sites.stat.washington.edu/people/raftery/Research/PDF/Gneiting2007jasa.pdf}

\bibitem[Gneiting y Ranjan(2011)]{gneiting2011}
Gneiting, T. y Ranjan, R. (2011). Comparing density forecasts using threshold- and quantile-weighted
scoring rules. \emph{Journal of Business \& Economic Statistics}, 29(3), 411--422.

\bibitem[Kalman(1960)]{kalman1960}
Kalman, R.~E. (1960). A new approach to linear filtering and prediction problems. \emph{Journal of Basic
Engineering}, 82(1), 35--45.

\bibitem[Koenker y Bassett(1978)]{koenker1978}
Koenker, R. y Bassett, G. (1978). Regression quantiles. \emph{Econometrica}, 46(1), 33--50.

\bibitem[K{\"u}nsch(1989)]{kunsch1989}
Künsch, H.~R. (1989). The jackknife and the bootstrap for general stationary observations. \emph{Annals of
Statistics}, 17(3), 1217--1241.

\bibitem[Ledoit y Wolf(2004)]{ledoit2004}
Ledoit, O. y Wolf, M. (2004). A well-conditioned estimator for large-dimensional covariance matrices.
\emph{Journal of Multivariate Analysis}, 88(2), 365--411.

\bibitem[Lee(2004)]{lee2004}
Lee, R.~W. (2004). The moment formula for implied volatility at extreme strikes. \emph{Mathematical
Finance}, 14(3), 469--480.

\bibitem[L{\'o}pez~de~Prado(2018)]{lopezdeprado2018}
López de Prado, M. (2018). \emph{Advances in Financial Machine Learning}. Wiley.

\bibitem[Martin(2017)]{martin2017}
Martin, I. (2017). What is the expected return on the market? \emph{Quarterly Journal of Economics},
132(1), 367--433.

\bibitem[Martin y Wagner(2019)]{martin2019}
Martin, I. y Wagner, C. (2019). What is the expected return on a stock? \emph{Journal of Finance}, 74(4),
1887--1929.

\bibitem[Mohajerin~Esfahani y Kuhn(2018)]{esfahani2018}
Mohajerin Esfahani, P. y Kuhn, D. (2018). Data-driven distributionally robust optimization using the
Wasserstein metric: performance guarantees and tractable reformulations. \emph{Mathematical Programming},
171, 115--166. arXiv:1505.05116. \url{https://arxiv.org/abs/1505.05116}

\bibitem[Patton(2011)]{patton2011}
Patton, A.~J. (2011). Volatility forecast comparison using imperfect volatility proxies. \emph{Journal of
Econometrics}, 160(1), 246--256.

\bibitem[Petersen y M{\"u}ller(2016)]{petersen2016}
Petersen, A. y Müller, H.-G. (2016). Functional data analysis for density functions by transformation to a
Hilbert space. \emph{Annals of Statistics}, 44(1), 183--218.

\bibitem[Peyr{\'e} y Cuturi(2019)]{peyre2019}
Peyré, G. y Cuturi, M. (2019). Computational optimal transport. \emph{Foundations and Trends in Machine
Learning}, 11(5--6), 355--607. arXiv:1803.00567. \url{https://arxiv.org/abs/1803.00567}

\bibitem[Rockafellar y Uryasev(2000)]{rockafellar2000}
Rockafellar, R.~T. y Uryasev, S. (2000). Optimization of conditional value-at-risk. \emph{Journal of
Risk}, 2(3), 21--41.

\bibitem[Ross(2015)]{ross2015}
Ross, S. (2015). The recovery theorem. \emph{Journal of Finance}, 70(2), 615--648.

\bibitem[Timmermann(2006)]{timmermann2006}
Timmermann, A. (2006). Forecast combinations. En \emph{Handbook of Economic Forecasting}, vol.~1,
135--196. Elsevier.

\bibitem[Zhang et~al.(2022)]{zhang2022}
Zhang, C., Kokoszka, P. y Petersen, A. (2022). Wasserstein autoregressive models for density time series.
\emph{Journal of Time Series Analysis}, 43(1), 30--52. arXiv:2006.12640. \url{https://arxiv.org/abs/2006.12640}

\end{thebibliography}
